% R038：例2的统计量计算流程；箭头只连接相邻框。
\begin{center}
	\begin{tikzpicture}[
			x=1cm,
			y=1cm,
			>=Stealth,
			every node/.style={font=\small},
			box/.style={
					draw=CExample,
					fill=green!4,
					rounded corners=5pt,
					align=center,
					font=\scriptsize,
					inner sep=4pt
				},
			givenbox/.style={
					box,
					minimum width=2.5cm,
					minimum height=1.55cm
				},
			meanbox/.style={
					box,
					minimum width=2.25cm,
					minimum height=1.05cm
				},
			calcbox/.style={
					box,
					minimum width=3.15cm,
					minimum height=1.08cm
				},
			resultbox/.style={
					box,
					minimum width=2.75cm,
					minimum height=1.55cm
				},
			arr/.style={
					->,
					CExample,
					very thick,
					rounded corners=3pt
				}
		]

		\path[use as bounding box] (0,0) rectangle (12.4,5.6);

		\node[
			font=\large\bfseries,
			text=CExample
		] at (6.2,5.23)
		{例2：统计量按“均值—离差和—回归系数”顺序计算};

		%==================== 第一列：已知统计量 ====================

		\node[givenbox] (raw) at (1.40,3.25)
		{
			$n=5$\\[-2pt]
			$\sum x_i=15,\quad \sum y_i=20$\\[-2pt]
			$\sum x_iy_i=70$\\[-2pt]
			$\sum x_i^2=55$
		};

		%==================== 第二列：样本均值 ====================

		\node[meanbox] (mean) at (4.10,3.25)
		{
			$\bar x=15/5=3$\\
			$\bar y=20/5=4$
		};

		%==================== 第三列：两条独立支路 ====================

		\node[calcbox] (sxy) at (7.25,4.08)
		{
			$S_{xy}=\sum x_iy_i-n\bar x\bar y$\\
			$=70-5\cdot3\cdot4=10$
		};

		\node[calcbox] (sxx) at (7.25,2.42)
		{
			$S_{xx}=\sum x_i^2-n\bar x^2$\\
			$=55-5\cdot3^2=10$
		};

		%==================== 第四列：回归方程 ====================

		\node[resultbox] (final) at (10.70,3.25)
		{
			$\hat b=\dfrac{S_{xy}}{S_{xx}}=1$\\[2pt]
			$\hat a=\bar y-\hat b\bar x=1$\\[2pt]
			$\hat y=x+1$
		};

		%==================== 箭头 ====================

		% 已知统计量 → 均值
		\draw[arr]
		(raw.east)--(mean.west);

		% 均值 → Sxy：先水平，再向上，最后水平进入
		\draw[arr]
		([yshift=1.8mm]mean.east)
		-- ++(.24,0)
		|- (sxy.west);

		% 均值 → Sxx：先水平，再向下，最后水平进入
		\draw[arr]
		([yshift=-1.8mm]mean.east)
		-- ++(.24,0)
		|- (sxx.west);

		% Sxy → 结果框上部
		\draw[arr]
		(sxy.east)
		-- ++(.24,0)
		|- ([yshift=4.2mm]final.west);

		% Sxx → 结果框下部
		\draw[arr]
		(sxx.east)
		-- ++(.24,0)
		|- ([yshift=-4.2mm]final.west);

		%==================== 底部说明 ====================

		\node[
			font=\scriptsize,
			text=black!65,
			text width=11.4cm,
			align=center
		] at (6.2,.62)
		{
			两条支路分别计算离差乘积和与 $x$ 的离差平方和；\\[-1pt]
			再合并得到 $\hat b$，最后利用样本中心求出 $\hat a$。
		};

	\end{tikzpicture}
\end{center}